\documentclass{article}
\usepackage[top=2cm,bottom=2cm,left=2cm,right=2cm]{geometry}
\usepackage{amsmath, amssymb, amsthm}
\usepackage{hyperref}
\usepackage{csquotes}

\usepackage{datetime}
\newdateformat{petsa}{\the\day\ \shortmonthname[\the\month] \the\year}

\newtheorem{definition}{Definition}[section]
\newtheorem{theorem}{Theorem}[section]
\newtheorem{lemma}[theorem]{Lemma}
\newtheorem{problem}[theorem]{Problem}

\title{Either This Sentence is True or $1+1=2$}
\author{poypoyan}
\date{\petsa\today}   % \petsa\today

\begin{document}

\maketitle

\begin{abstract}
This is an analysis of Liar and Curry's paradox, and through this we get variants, such as the title above.
\end{abstract}

\section{The Liar and Curry's Paradox}
We work in classical logic. The Liar paradox:
\begin{displayquote}
This sentence is false.
\end{displayquote}
and Curry's paradox:
\begin{displayquote}
If this sentence is true, then $1=2$.
\end{displayquote}
have two properties in common (at least on how we presented these two): 1) both are self-referential sentences, and 2) both purely deal with truth and falsity of sentences. There is an assumption on sentences that \textit{must always hold}:
\begin{align}\label{main-assump}
A \iff A
\end{align}
With the self-reference, this can be interpreted as \textit{the equivalence of the truth value of ``this sentence" and the truth value of the actual content of the sentence}. Now we can formalize the paradoxes by substituting the formalized content to the right hand side. For Liar paradox:
\begin{align*}
&A \iff (A \rightarrow \bot) \\
\implies &A \iff \neg A
\end{align*}
It is clear that \ref{main-assump} is \textit{false/contradicted} here. Now let's formalize Curry's paradox:
\[A \iff A \rightarrow B\]
If $B$ is false, then we are back to the formalized Liar paradox! On the other hand, in the ``classical" analysis (e.g., in \cite[An Informal Argument]{sep-curry-paradox}), the sentence makes us conclude that $B$ must be true. The convergence of the two analyses is on the fact that with the falsity of $B$, there is now contradiction. The difference is on what is involved in the contradiction: $B$ itself in the classical analysis, in contrast to \ref{main-assump} in our analysis.

\section{Synthesis}
The idea is that making \ref{main-assump} false can indicate the ``paradoxness" of the two paradoxes. Can we find variants that satisfy the two aforementioned properties? Let's stick on one independent proposition $B$:
\[A \iff P(A, B)\]
What about the truth of $B$ causes the contradiction instead? One such $P$ would be $\neg A \vee \neg B \iff \neg (A \wedge B)$ (setting $B$ to true simplifies to $\neg A$). Translating back to less formal, this is
\begin{displayquote}
This sentence and ``$1+1=2$", cannot be both true, or\\
It's not the case that this sentence is true and $1+1=2$.
\end{displayquote}
Here is the corresponding classical analysis: Suppose the sentence is false. Then both ``this sentence" and ``$1+1=2$" are true, hence the sentence is actually true. Now suppose the sentence is true. Then ``$1+1=2$" must be false. That's not right.

Another such $P$ would be $(\neg A \vee \neg B ) \wedge (A \vee B) \iff A \veebar B$, which translates back to the title of this note:
\begin{displayquote}
Either this sentence is true or $1+1=2$.
\end{displayquote}
Here is the corresponding classical analysis: Suppose the sentence is false. Then either the two are both true or both false. As we showed in the previous classical analysis, both true implies that the sentence is actually true. This leaves us with both false, which implies that ``$1+1=2$" must be false. Now suppose the sentence is true. This also implies that ``$1+1=2$" must be false.

I find this one more aesthetically pleasing than the former because it (in English) looks similar to Curry's paradox, wherein $B$ is not under quotation and there is no direct mention of ``not" or ``false".

\section{Closing Remarks}
This analysis can be extended to multiple independent propositions or to multiple sentences referencing each other. Moreover, other paradoxes can be subject to this analysis by making explicit the assumption/s that ``must always hold" and that the paradox contradicts. In some sense, these assumptions are contexts that put the paradoxes/contradictions ``under containment", which leads us to a further elucidation of paraconsistent reasoning.

\bibliographystyle{plain}
\bibliography{liarcurry}
\end{document}
